% ============================================================
% Page geometry and chapter style
% ============================================================

% Page geometry: B5-like trimmed page (250 mm x 176 mm).
\stockaiii
\pageaiv
\setstocksize{250mm}{176mm}
\settrims{0pt}{0pt}
\settrimmedsize{\stockheight}{\stockwidth}{*}
\setulmargins{2.5cm}{*}{*}
\setlrmarginsandblock{2cm}{1.5cm}{*}
\checkandfixthelayout

% Paragraphing. With memoir's default \parskip of 0pt the only mark of a
% paragraph break is the first-line indent, so a long stretch of prose with
% no display or figure in it sets as an undifferentiated slab. A fraction
% of a line between paragraphs breaks that up while staying well short of
% the block-paragraph look; the indent is kept, since it is what still
% distinguishes a new paragraph from a line that happens to end short. Set
% through memoir's \abnormalparskip rather than \setlength directly, so
% list and display spacing follow \parskip instead of drifting out of
% proportion with it.
\abnormalparskip{0.4\onelineskip plus 0.15\onelineskip minus 0.05\onelineskip}

% Page style and chapter headings

% memoir's `ruled` marks put two spaces after the number, since they are
% built with `{. \space}` — visible now that the section mark runs in the
% header ("1.1.  Title ..."). Re-issue them with a single space, appended
% to \@ruledmarks so it still applies whenever the style is selected.
\makeatletter
\g@addto@macro\@ruledmarks{%
	\createmark{chapter}{left}{shownumber}{}{.\space}%
	\createmark{section}{right}{shownumber}{}{.\space}%
}
\makeatother

\pagestyle{ruled}
\chapterstyle{madsen}
% Use Roboto for chapter/appendix heading text.
\renewcommand*{\chapnamefont}{\roboto\Large\scshape\raggedleft}
\renewcommand*{\chaptitlefont}{\roboto\Huge\bfseries\raggedleft}
% Keep madsen style but align Chapter+number to the text block.
% The default madsen uses a 0pt box that can push the number into the margin.
\makeatletter
\renewcommand*{\printchaptername}{}
\renewcommand*{\chapternamenum}{}
% The numeral alone carries the red accent at full brand saturation — one
% deliberate hit of institutional colour per chapter opener, structural
% rather than inline, so it never competes with the darkened link colours
% in the running text (see preamble/colors.tex).
\renewcommand*{\printchapternum}{%
    \makebox[\textwidth][r]{%
        \resizebox{!}{4ex}{\chapnamefont\bfseries\sffamily\@chapapp\ \textcolor{accent_red}{\thechapter}}%
    }%
}
\renewcommand*{\printchapternonum}{%
    \afterchapternum
}
\makeatother
\setlength{\beforechapskip}{15pt}

% Section numbering scheme (memoir defaults, stated explicitly):
% - Chapters:       1, 2, ...
% - Sections:       1.1, 1.2, ... (chapter-qualified)
% - Subsections:    1.1.1, 1.1.2, ...
% - Subsubsections: 1.1.1.1, 1.1.1.2, ...
\renewcommand*{\thechapter}{\arabic{chapter}}
\renewcommand*{\thesection}{\thechapter.\arabic{section}}
\renewcommand*{\thesubsection}{\thesection.\arabic{subsection}}
\renewcommand*{\thesubsubsection}{\thesubsection.\arabic{subsubsection}}
\setsecnumdepth{subsubsection}
\settocdepth{subsection}

% Use Roboto for sectioning titles.
\setsecheadstyle{\roboto\Large\bfseries}
\setsubsecheadstyle{\roboto\large\bfseries}
\setsubsubsecheadstyle{\roboto\normalsize\bfseries}
\setbeforeparaskip{1.0\onelineskip plus 0.5\onelineskip minus 0.2\onelineskip}
\setafterparaskip{0.5\onelineskip}
\setparaheadstyle{\robotoMediumFace\normalsize\itshape}
\setsubparaheadstyle{\robotoMediumFace\normalsize}

% Prevent a \section heading from being stranded near the bottom of a
% page with almost no body text below it before the next page break.
\pretocmd{\section}{\needspace{6\baselineskip}}{}{}

% Running heads follow the TU Delft dissertation layout: the page number
% sits in the header at the *outer* edge (left on versos, right on rectos)
% and the running mark at the inner edge, so a spread reads number-outward
% on both sides.  The footer is empty; the only page carrying a number at
% the foot is a chapter opening, centred, through the plain style below.
%
% memoir's own `ruled` marks are already what we want and are left alone:
% \createmark{chapter}{left} puts the chapter in \leftmark *and clears*
% \rightmark, and \createmark{section}{right} then fills \rightmark — so a
% recto before a chapter's first section, or in a chapter with no sections
% at all, simply has a blank mark.
\makeevenhead{ruled}{\roboto\thepage}{}{\roboto\leftmark}
\makeoddhead{ruled}{\roboto\rightmark}{}{\roboto\thepage}
\makeevenfoot{ruled}{}{}{}
\makeoddfoot{ruled}{}{}{}

% Chapter-opening pages use memoir's plain pagestyle (aliased to `chapter`);
% this is the one place the number stays at the foot, centred.
\makeevenfoot{plain}{}{\roboto\thepage}{}
\makeoddfoot{plain}{}{\roboto\thepage}{}

% Graphics search path. Inside a subfile, \includegraphics also resolves
% paths relative to that subfile's own directory (subfiles package), so
% chapters can write `figures/foo.pdf` both in the full build and in a
% standalone chapter build.
\graphicspath{
    {frontmatter/img/}
    {appendix/img/}
}
